% =========================================================
% Proof: Euclid's Lemma (Coprimality Form)
% Source: volume-iii/algebra/abstract-algebra/notes/notes-integers.tex
% =========================================================

\subsection*{Euclid's Lemma (Coprimality Form)}
\label{prf:euclids-lemma-coprimality}

\begin{remark*}[Return]
\hyperref[cor:euclids-lemma-coprimality]{$\leftarrow$ Back to Corollary
(Euclid's Lemma --- Coprimality Form) in Notes}
\end{remark*}

\begin{proof}
\Claim Let $a, b, c \in \mathbb{Z}$. If $a \mid bc$ and $\gcd(a,b) = 1$,
then $a \mid c$.

\Given $a \mid bc$ and $\gcd(a,b) = 1$.
\Goal To show $a \mid c$.

\ByThm{Bézout's Lemma}, since $\gcd(a,b) = 1$, there exist $m, n \in \mathbb{Z}$
such that
\[
    1 = ma + nb.
\]
Multiplying both sides by $c$:
\[
    c = mac + nbc.
\]
Since $a \mid a$, we have $a \mid mac$.
Since $a \mid bc$ by hypothesis, we have $a \mid nbc$.
\ByThm{Divisibility of Linear Combinations}, $a \mid mac + nbc = c$.

\Hence $a \mid c$. \AsReq
\end{proof}

\begin{remark*}[Proof shape]
Direct proof via Bézout's Lemma. The key move is multiplying the
Bézout relation $1 = ma + nb$ through by $c$, which writes $c$ itself
as a linear combination $mac + nbc$. The two hypotheses — $a \mid a$
(trivially) and $a \mid bc$ (given) — then each handle one term,
and divisibility of linear combinations closes the argument.
This is the standard ``Bézout trick'': turn a coprimality condition
into a linear combination, then scale to reach the desired conclusion.

\FlashProof{Direct proof via Bézout's Lemma. The key move is multiplying the Bézout relation $1 = ma + nb$ through by $c$, which writes $c$ itself as a linear combination $mac + nbc$. The two hypotheses — $a \mid a$ (trivially) and $a \mid bc$ (given) — then each handle one term, and divisibility of linear combinations closes the argument. This is the standard ``Bézout trick'': turn a coprimality condition into a linear combination, then scale to reach the desired conclusion.}
\FlashUses{\begin{itemize} \item \hyperref[cor:bezout-lemma]{Bézout's Lemma} — to obtain the linear combination $1 = ma + nb$ from $\gcd(a,b) = 1$. \item \hyperref[lem:div-linear-combo]{Lemma (Divisibility of Linear Combinations)} — since $a$ divides both terms, it divides their sum $c$. \end{itemize}}
\end{remark*}

\begin{remark*}[Dependencies]
\begin{itemize}
    \item \hyperref[cor:bezout-lemma]{Bézout's Lemma} — to obtain the
          linear combination $1 = ma + nb$ from $\gcd(a,b) = 1$.
    \item \hyperref[lem:div-linear-combo]{Lemma (Divisibility of Linear
          Combinations)} — since $a$ divides both terms, it divides
          their sum $c$.
\end{itemize}
\end{remark*}